\section{Reproducibility \& colophon}\label{app:repro}

This whitepaper is itself a small, self-contained build: a contributor who clones the
repository can regenerate the PDF, with every diagram, from the sources checked into
\srcf{docs/DesignWhitepaperLatex/}. The layout and style contract are inherited from the
Architecture Review (\srcf{docs/ArchitectureReviewLatex/}), which in turn inherited them
from the Part Placement design whitepaper (\srcf{docs/PartPlacementDesignLatex/}), so the
three documents read as one series. This appendix records the build pipeline, the
revision every claim is pinned to, and how the document was produced and verified.

\subsection{Directory layout and pipeline}

The master file \srcf{DesignWhitepaper.tex} carries only the document scaffolding ---
title block, front matter, and an ordered list of \cppinline|\input| directives --- while
the prose lives in \srcf{sections/}, one file per section or appendix. The shared style
contract is isolated in \srcf{preamble.tex}: the color palette, the fonts (XCharter text,
\code{newtxmath}, Inconsolata mono), the callout environments (including the
\code{physbox} added for this document, which keeps physics background visually separate
from statements about the code), and the domain macros (\code{\textbackslash src},
\code{\textbackslash srcf}, \code{\textbackslash tikzfig}, the feature-status markers)
that the section files depend on. A section file is just a body, never a standalone
\LaTeX{} program.

Every figure is hand-built Ti\emph{k}Z under \srcf{diagrams/tikz/}, included directly by
\code{\textbackslash tikzfig} and compiled in line by the main engine --- there is no
PlantUML stage and no external image asset. The UML class, sequence, and component
figures use dedicated styles added to \srcf{diagrams/tikzstyles.tex} for this document
(three-compartment class boxes, inheritance/realization/composition arrowheads,
lifelines). The build is two commands:

\begin{minted}{text}
make            # latexmk -pdf -shell-escape DesignWhitepaper.tex
make veryclean  # remove aux files, the minted cache, and the PDF
\end{minted}

The engine is \code{pdflatex}, driven by \code{latexmk} so the multi-pass dance
(cross-references, the table of contents, \code{cleveref}, the bookmark tree) converges
automatically. The one non-default requirement is \code{-shell-escape}: listings are
typeset by \code{minted}, which shells out to \code{pygmentize}. Prerequisites are a
\TeX{} distribution providing \code{latexmk}, \code{pdflatex}, and \code{minted} with its
\code{pygmentize} helper.

\subsection{The revision under description}

The document describes one tree state: commit \code{254cd41} (``Major comment
cleanup''), the tip of the \code{development} line at the time of writing, described
against a byte-identical working tree in July 2026. Every \code{\textbackslash src}
citation refers to that revision; line numbers will drift as the tree moves, so a reader
on a later commit should treat them as anchors into \code{254cd41}, not absolutes.

The ground-truth policy is the one stated in \cref{sec:intro}: the source at the pinned
commit and the two companion documents --- the Architecture Review and the Part Placement
design whitepaper --- are authoritative; every other artifact under \srcf{docs/} is
historical and was deliberately not relied on. Where a companion document and the code
disagree, this document reports the code and says so in place.

\subsection{How this document was produced and verified}

The document was produced in three staged passes, each a fan-out of independent
workers.

\emph{Deep reads.} Twelve parallel deep-reads produced working dossiers: nine full reads
of the code subsystems (the part tree and placement machinery, the controller and
\code{utils} foundation, motors, the simulation core, the environment models,
aerodynamics, the GUI and visualizer, the CLI and persistence, and the build/test/CI
stack), a digest of the two companion documents, and two research dossiers on OpenRocket
--- architecture and simulation methods --- compiled from primary sources
(\texttt{openrocket.info}, the OpenRocket documentation, Niskanen's thesis and the
technical documentation, and the \texttt{openrocket/openrocket} repository at
\code{release-24.12} and \code{unstable}, fetched July 2026). Each dossier carried a
citation ledger of exact \code{file:line} anchors re-checked against the source before
use.

\emph{Writing.} Each section was then written from its dossiers --- writers were
forbidden to assert anything about the code not present in a dossier or re-verified in
the source directly, with general physics theory admitted only as clearly-framed
background. The diagrams were authored the same way and test-compiled standalone before
inclusion.

\emph{Adversarial verification.} Finally, one verifier per section re-checked the
finished text against the source itself: every citation opened and its line number and
substance confirmed, every listing diffed verbatim, every equation re-derived against
the code, every count and constant re-taken (over two thousand checks; roughly eighty
line-number corrections and seventy substantive claim corrections were applied). Three
cross-cutting passes followed: a consistency sweep that adjudicated shared facts across
sections against the source, a diagram pass that re-verified twelve of the thirteen
figures element-by-element and recompiled them (the thirteenth, the front-end component
map, was authored afterward from already-verified section text and test-compiled the
same way), and a completeness critique whose findings were folded back into the text. OpenRocket-side claims were verified for
faithful transmission from the research dossiers; the few that could not be pinned to a
release tag are marked as branch-verified where they appear.

Two behaviors were confirmed by execution during the underlying review series rather
than by reading: the CLI command inventory (\cref{sec:interfaces}) and end-to-end flight
behavior through the built \code{qtrocket-cli} binary. The dossiers themselves are
working notes, not part of the document, and live untracked under \code{notes/}.
